\section{Soundness over unrestricted integers}\label{sec:soundness}

The proof first excludes a zero multiplier. It then applies the valuation
lemma before making any positivity assumption about the guard. This order
is needed because all twelve coordinates are unrestricted integers.

\subsection{The zero-multiplier boundary}

\begin{lemma}\label{lem:nonzero-scale}
Neither \eqref{eq:primary} nor \eqref{eq:companion} has an integer zero
with $S=0$. At every zero, $I,s,r,c$ are nonzero and $Q>1$.
\end{lemma}
\begin{proof}
By Lemma~\ref{lem:residues}, $H\ne0$. Thus $S=0$ implies $I=0$,
$Q=0$, and $\delta=-1$. All scaled radicands except $B^2\Dm$ vanish,
and both final polynomials specialize to
\[
 \bigl(A^2-B^2\Dm\bigr)^{16}.
\]
The two factors of the inner expression over $\Q(\sqrt{\Dm})$ are
\[
 A\pm B\sqrt{\Dm}
 =H^5\pi\!\left(\frac{u\pm\sqrt{\Dm}}{H}\right).
\]
Since $\Dm\ge1$, both arguments belong to a real extension of $\Q$ of
degree at most two. Lemma~\ref{lem:intersective} says neither factor can
vanish. Therefore $S\ne0$. Now $I=src^2$ is nonzero, so each of
$s,r,c$ is nonzero; $Q=(\Delta I)^2>1$ follows from $\Delta>1$.
\end{proof}

\subsection{Root extraction without a sign assumption}

Suppose first that $\primary=0$. The unconditional fact $\delta\ne0$
allows the rational interpretation in \eqref{eq:primary}. In a splitting
field, one factor of the norm vanishes. After division by $S\ne0$,
the five radical terms, up to independent signs, are
\begin{equation}\label{eq:normalized-roots}
 \frac{B\sqrt{\Dm}}{S},\qquad
 \sqrt{\Cc/\delta},\qquad
 3\sqrt{\Df},\qquad9\sqrt{\Dt},\qquad27\sqrt{\Dg}.
\end{equation}
Their signed sum is the rational number $n-A/S$. The first four terms
are nonzero by the residues in Lemma~\ref{lem:residues}. Their squared
$3$-adic orders are, respectively,
\begin{equation}\label{eq:orders}
 -2\vt(S),\qquad0,\qquad2,\qquad4.
\end{equation}
Here $\vt(S)=\vt(I)+5\vt(H)\ge5$. If $\Dg\ne0$, the last
squared order is
\[
 6+\vt(\Dg)\ge6;
\]
it need not be even. If $\Dg=0$, omit the zero term. Every subset
with at least two surviving terms has a unique even minimum, because
any such subset contains one of the first four terms. Thus
Lemma~\ref{lem:separation} makes all terms in
\eqref{eq:normalized-roots} rational. This argument also covers
$\Dg<0$; no realness premise was used.

Since $B/S\ne0$, it follows that $\Dm,\Df,\Dt,\Dg$ have rational
square roots. A rational square root of an integer is an integer, so
there are integers $d,f,\tau,g$ with
\begin{equation}\label{eq:extracted-integer-roots}
 d^2=\Dm,\qquad f^2=\Df,\qquad
 \tau^2=\Dt,\qquad g^2=\Dg.
\end{equation}
The auxiliary root is initially rational; Section~\ref{sec:aux-integral}
will prove it integral.

For $\companion=0$, replace the second term of
\eqref{eq:normalized-roots} by $\sqrt{\delta\Cc}$. It again has
squared $3$-adic order zero, so the same proof supplies an integer
$\widetilde y$ with $\widetilde y^2=\delta\Cc$ and all four roots
in \eqref{eq:extracted-integer-roots}.

\subsection{The strict guards}

The extracted square $g^2=\Dg$ proves $\Dg\ge0$. Since $s$ is a
nonzero integer, $P\le0$ would give
\[
 \Dg=(3P-4)s^2+1\le-3.
\]
Hence $P>0$. We record the elementary guard principle to make this
sign inference explicit.

\begin{lemma}[Joint positivity]\label{lem:joint-positive}
For integers $A_1,A_2$, there exists an integer $t$ such that
\[
 A_1A_2-(t^2-A_1-A_2+1)^2>0
\]
if and only if $A_1>0$ and $A_2>0$.
\end{lemma}
\begin{proof}
The inequality implies $A_1A_2>0$. If both factors were negative, then
$t^2-A_1-A_2+1>|A_1|+|A_2|\ge2\sqrt{A_1A_2}$, a contradiction.
Conversely, put $M=A_1+A_2-1\ge1$. An integer square nearest to the
integer $M$ differs from $M$ by an integer of absolute value strictly
less than $\sqrt M$: between consecutive squares $a^2,(a+1)^2$, the
nearest-square error is at most $a$, and at an endpoint it is zero.
Since $M\le A_1A_2$, this choice proves the strict inequality.
\end{proof}

Applying the forward implication to $A_1=G$, $A_2=\eta\zeta$ yields
\begin{equation}\label{eq:strict-guards}
 G>0,\qquad \eta\zeta>0.
\end{equation}
The determinant guard also restores the signs required by the inherited
signed kernel. Set $U=\Acone F$, $W=\Bcone C$, and
$\alpha=q-F-C$. From $UW>(q-U-W)^2$ one obtains $UW>0$. If $q$
had the opposite sign to $U,W$, the square would be at least
$(|U|+|W|)^2\ge4UW$, impossible. Moreover
\[
 |q|>|U|+|W|-\sqrt{UW}\ge\tfrac12(|U|+|W|).
\]
Since both cone scales are at least two, $F,C,\alpha$ have the same
nonzero sign as $q$.

For clarity, $R>0$ already follows at this point. The authentic fixed
tuple has $0<M_C<\beta<M_F$. If $q>0$, then $J>0$ and
\[
 q(q-F)-C=q\alpha+(q-1)C>0.
\]
Both summands in the definition of $R$ are positive. If $q=-q_0<0$,
write $C=-C_0$, $\alpha=-\alpha_0$, with $C_0,\alpha_0>0$.
Then
\[
 q(q-F)-C=q_0(C_0+\alpha_0)+C_0>0,
 \quad J<0,\quad M_C-q_0M_F<0.
\]
Again both summands defining $R$ are positive. Thus no later rank
argument is being used to justify this preliminary sign fact.

\subsection{The rational auxiliary root is integral}\label{sec:aux-integral}

\begin{lemma}[Adjacent-coefficient descent]\label{lem:adjacent-descent}
Let $Q>1$ be an integer and let $V\in\Z$, $y\in\Q$ satisfy
\[
 QV^2-(Q-1)y^2=1.
\]
Then $y\in\Z$.
\end{lemma}
\begin{proof}
Changing signs makes $V,y$ positive. The norm excludes $V=0$.
Write $y=u/b$ in lowest terms with $b>0$. The equation implies
$b^2\mid Q-1$, so $(Q-1)y$ is integral. If $V>1$, use
\begin{align*}
 V'&=(2Q-1)V-2(Q-1)y,\\
 y'&=(2Q-1)y-2QV.
\end{align*}
This transform preserves the norm, and $V'$ is integral. Because
$2Q-1\equiv1\pmod b$, the reduced denominator of $y'$ is still
$b$. The norm gives $y>V$ and $(Q-1)y<QV$, hence
$-V<V'<V$. The norm also excludes $V'=0$. Taking absolute values
therefore gives descent in the positive integer $V$ without changing
$b$. At $V=1$ the equation gives $|y|=1$, so $b=1$.
\end{proof}

For the primary equation, the rational root $y$ extracted from
$\Cc/\delta$ satisfies the lemma's equation because
$\delta=Q-1$ and $\Cc=QV^2-1$. For the companion, take
$y=\widetilde y/\delta$; its square is again $\Cc/\delta$.
Lemma~\ref{lem:nonzero-scale} supplies $Q>1$ in both cases. Thus the
same integer auxiliary Pell witness is recovered from both polynomials.

\subsection{Recovering the main congruence}

Absorb the signs of a vanishing conjugate factor into the integral
roots. The primary equation gives
\begin{equation}\label{eq:root-sum}
 Sn=A+Bd+S(y+3f+9\tau+27g).
\end{equation}
For the companion, replace $y$ on the right by $\widetilde y$.
Consequently $H^5\mid A+Bd$ in both cases. Since
\[
 A+Bd=\Phi_H(HT+d-b_0),
\]
this restores exactly the missing congruence. Indeed, if $p^e\Vert H$
and $v_p(HT+d-b_0)=b<e$, the leading fifth power in
\[
 \Phi_H(z)=z^5+3H^2z^3-19H^3z^2-57H^5
\]
has valuation $5b$, strictly below all three other terms. Divisibility
by $p^{5e}$ is impossible. Doing this for each prime divisor of $H$
proves $H\mid HT+d-b_0$, and therefore
\begin{equation}\label{eq:gamma-restored}
 \gamma=\frac{d-b_0}{H}\in\Z,\qquad d=X+ac+\gamma H.
\end{equation}
This introduces no independent coordinate; $\gamma$ is reconstructed
from an existing zero.

\subsection{The precise inherited signed-kernel interface}

The recovered equations are
\begin{align}
 d^2-\Delta c^2&=1,& d&=X+ac+\gamma H,\nonumber\\
 f^2-\Delta I^2&=1,& I&=src^2\ne0,\nonumber\\
 (\Delta I)^2V^2-((\Delta I)^2-1)y^2&=1,&
 V&=f^2(c-R+4IT)-c,\nonumber\\
 \tau^2-XY^2(XY^2+1)k^2&=1,&
 G&>0,\quad\eta\zeta>0. \label{eq:restored-kernel}
\end{align}
The generalized signed-kernel proof used by the incoming twelve-variable
construction is preserved as
\path{integer12/proofs/INTERPOLATION12.md}, Sections 5--7, and its
independent arithmetic audit as
\path{integer12/proofs/INDEPENDENT_TWELVE_AUDIT.md}, Sections 2--4.
They specify which ordinary rank, orientation, and native-decoding
theorems are inherited. We give the normalized strong-rank argument here
because it prevents an incorrect parity shortcut.

For an integer $A_*>1$, define the Pell polynomials by
\begin{equation}\label{eq:pell-def}
 \chi_{A_*}(m)+\psi_{A_*}(m)\sqrt{A_*^2-1}
   =(A_*+\sqrt{A_*^2-1})^m,\qquad m\ge0.
\end{equation}
The inherited first/main rank bootstrap, applied to the first and last
Pell equations and the strict guards in \eqref{eq:restored-kernel},
gives
\[
 c_*=|c|=\psi_{A_*}(p),\qquad A_*=|a+2|,\qquad p\ge25.
\]
The normalized strong equation in \eqref{eq:restored-kernel} gives
\[
 |I|=\psi_{A_*}(m),\qquad |f|=\chi_{A_*}(m),\qquad m\ge1.
\]
This uses the usual integral Pell classification for
$A_*^2-1$: $A_*+\sqrt{A_*^2-1}$ is the fundamental positive unit
with integer coefficients, as follows by descent on the first coordinate.

Since $c_*^2\mid |I|$, strong divisibility of the Lucas sequence
$\psi_{A_*}$ gives $p\mid m$. Write $m=p\ell$. Expanding the
$\ell$th power of the unit at index $p$ gives
\[
 \frac{\psi_{A_*}(p\ell)}{c_*}
 \equiv\ell\chi_{A_*}(p)^{\ell-1}\pmod{c_*}.
\]
The left side is divisible by $c_*$, and the Pell identity gives
$\gcd(\chi_{A_*}(p),c_*)=1$. Hence $c_*\mid\ell$ and, in
particular, $pc_*\mid m$. These are the strong-rank inputs needed
by the inherited signed auxiliary descent. They do not require $I/c^2$
to be odd.

For the auxiliary interface put $\rho=\Delta I$. The discarded
exposed quotient coordinates are restored by the integer polynomials
\begin{align*}
 o&=f(c-R+4IT),\\
 j_{\rm aux}&=\Delta s^2r^2c^3(c-R)+4src\Df T.
\end{align*}
Direct expansion gives
\[
 V=fo-c=cj_{\rm aux}-R,
 \qquad \rho^2=\Delta(f^2-1).
\]
Also $\rho^2\equiv1-A_*^2\pmod{|f|}$ and
$\rho^2\equiv0\pmod{c_*}$. In the auxiliary Pell classification,
the first coordinate is $|\rho V|$. Its index is odd because
\[
 \chi_{|\rho|}(2j)\equiv(-1)^j\pmod{|\rho|},
 \qquad |\rho|>1,
\]
whereas $\rho\mid\rho V$. This step likewise needs no oddness of
$\rho$.

The inherited signed auxiliary step-down now yields $p=R$ and
$2n_{\rm first}=R+1$; its main-root residue gives
$X=\operatorname{sgn}(c)2^R$. In particular $R$ is odd, so the
fixed residue $X\equiv2\pmod3$ selects $c>0$ and $X=2^R$.
The inherited first-index sign identity gives $k>0$, and
$c/k=Y+\eta/k$, with $0<\eta/k<1$ and $|Y|>1$, then gives
$Y>0$. Finally $q\mid2^R$, $q\equiv1\pmod\beta$, and the
inherited Mersenne-radix sign lemma give $q>0$. The post-power native
decoder and its fixed blank-input calibration supply a finite syntactic
ZFC refutation. The relevant native proofs are retained under
\path{integer12/inherited/root_fusion/provenance/}, including
\texttt{PROOF44.md} and the \texttt{upstream\_ordinary/} calibration
documents. This is a composition of their ordinary theorems, rather than
a claim that finite circuit checks establish the unbounded compiler.

\begin{remark}
An attempted restoration $r_{\rm old}=(\Delta sr-1)/2$ would be
invalid when $\Delta sr$ is even. The normalized rank argument above
is the reason arbitrary nonzero multipliers are admissible. The proof
does not silently import the older odd-multiplier coordinate chart.
\end{remark}

\section{Completeness and the interpolated residue}\label{sec:completeness}

The only new completeness choice is a residue for the shorter quintic.
We reproduce the interpolation mechanism because its quantifier order
is essential: one strong Pell witness must be fixed before the requested
residue class is selected.

The inherited canonical native construction supplies the outer data and
main Pell roots with
\begin{equation}\label{eq:canonical}
 \begin{gathered}
 A_*=a+2>1\text{ even},\qquad R\equiv3\pmod4,\\
 c=\psi_{A_*}(R)>0,\qquad d=\chi_{A_*}(R),\qquad
 \gamma=(d-X-ac)/H\in\Z,\\
 G>0,\qquad\eta\zeta>0,\qquad \tau^2=\Dt.
 \end{gathered}
\end{equation}
In particular $c\equiv R\pmod4$, since the recurrence gives
$\psi_{A_*}(j+2)\equiv-\psi_{A_*}(j)\pmod4$ when $A_*$ is even.
The source-34 tuple's bounded recipe and the exact canonical hypotheses
are recorded in Appendix~\ref{app:coeff} and the inherited coefficient
proof. This is the canonical family used in the incoming interpolation
theorem, not a restriction imposed on arbitrary zeros.

\begin{lemma}[Full residue interpolation]\label{lem:interpolation}
Let $A_*>1$ be even, $\Delta=A_*^2-1$, $R\equiv3\pmod4$, and
$c=\psi_{A_*}(R)>0$. For every prescribed $N>0$, there is a fixed
positive pair $(I,f)$ such that
\[
 f^2-\Delta I^2=1,\qquad4Nc^2\mid I,
\]
and, for this same pair, every residue class $a_0\pmod I$ is represented
by an integer $T$ for which
\begin{equation}\label{eq:interpolation-equations}
 V=f^2(c-R+4IT)-c,\qquad
 (\Delta I)^2V^2-((\Delta I)^2-1)y^2=1
\end{equation}
has an integer solution $y$.
\end{lemma}
\begin{proof}
The Pell recurrence matrix has determinant one, so it has finite order
modulo $4Nc^2$. Choose a positive even multiple $m$ of that order and
put
\[
 I=\psi_{A_*}(m),\qquad f=\chi_{A_*}(m).
\]
Then $4Nc^2\mid I$, $I$ is even, $f$ is odd, and $\gcd(f,I)=1$.

We first justify the valuation fact
\begin{equation}\label{eq:lucas-valuation}
 v_p(I)\ge v_p(m)\qquad(p\mid I).
\end{equation}
For an odd prime dividing $\Delta$, the recurrence reduces to
$\psi_{A_*}(j)\equiv jA_*^{j-1}\pmod p$, so the first zero
index is $p$. For an odd prime not dividing $\Delta$, the norm-one
group in the split quadratic algebra or quadratic field over
$\mathbf F_p$ has order dividing $p-1$ or $p+1$, respectively; the
first zero index is therefore prime to $p$. If $p\mid\psi_{A_*}(j)$,
binomial expansion of the corresponding norm-one unit shows that
multiplying $j$ by $p$ raises the second-coordinate valuation by
exactly one, while multiplying it by a number prime to $p$ leaves
that valuation unchanged. In the first-zero case the initial valuation
is at least one, proving \eqref{eq:lucas-valuation} in both odd-prime
cases. At $p=2$, because $A_*$ is even and $m$ is even, the doubling
identity gives
\[
 v_2(\psi_{A_*}(m))=v_2(m)+v_2(A_*).
\]
This proves the remaining case, and hence $\gcd(m,I^2)\mid I$.

Write $\alpha=A_*+\sqrt\Delta$, so
$\alpha^m=f+I\sqrt\Delta$. In the integer-coefficient quadratic
algebra modulo $f^2$, its norm-one relation gives
\[
 \alpha^{4m}\equiv1-4f\alpha^m\pmod{f^2},
 \qquad \alpha^{4mf}\equiv1\pmod{f^2}.
\]
The second congruence follows by raising the first to the positive
power $f$. Thus $\lambda=4mf$ is a period of both Pell coordinates
modulo $f^2$. Since $\gcd(f,I)=1$ and
\eqref{eq:lucas-valuation} holds,
\[
 \gcd(\lambda,4I^2)\mid4I.
\]
For any specified class $a_0\pmod I$, the two congruences
\begin{equation}\label{eq:interpolation-crt}
 L\equiv R\pmod\lambda,\qquad
 L\equiv R-4Ia_0\pmod{4I^2}
\end{equation}
are therefore compatible. Choose a positive solution $L$. It is
$3\pmod4$. Put $\rho=\Delta I$ and
\[
 V_0=\chi_\rho(L)/\rho,\qquad y=\psi_\rho(L).
\]
The quotient $V_0$ is integral because the index is odd.

For completeness, the exact odd Chebyshev polynomials
$Q_j(X)=\chi_b(2j+1)/b$, $X=b^2$, satisfy
\[
 Q_0=1,\quad Q_1=4X-3,\quad
 Q_{j+1}=(4X-2)Q_j-Q_{j-1}.
\]
Induction gives
\[
 Q_j(X)\equiv(-1)^j(2j+1)\pmod{4X},\qquad
 Q_j(1-A_*^2)=(-1)^j\psi_{A_*}(2j+1).
\]
Since $\rho^2=\Delta(f^2-1)\equiv1-A_*^2\pmod{f^2}$,
these identities and \eqref{eq:interpolation-crt} imply
\begin{equation}\label{eq:two-aux-residues}
 V_0\equiv-c\pmod{f^2},\qquad
 V_0\equiv-R+4Ia_0\pmod{4I^2}.
\end{equation}
Here $L\equiv3\pmod4$ makes $(L-1)/2$ odd; the period $\lambda$
identifies $\psi_{A_*}(L)$ with $c$ modulo $f^2$.

Because $c-R$ is divisible by four and $f^2=1+\Delta I^2$,
the integer
\[
 N_0=V_0+c-f^2(c-R)
\]
is divisible by $f^2$ and congruent to $4Ia_0\pmod{4I^2}$.
The coprimality $\gcd(f^2,4I^2)=1$ now gives
\[
 T=\frac{N_0}{4If^2}\in\Z,\qquad T\equiv a_0\pmod I,
\]
where the last step also uses $f^2\equiv1\pmod I$.
Rearrangement supplies the first equation in
\eqref{eq:interpolation-equations}; the Pell identity for
$\chi_\rho(L),\psi_\rho(L)$ supplies the second. The initial pair
$(I,f)$ depended only on $N$ and the canonical data, not on $a_0$.
\end{proof}

\begin{remark}
The squared-modulus period is $4mf$, not $4m$. The latter is generally
incorrect. Also, even $I$ and odd $f$ are choices made for completeness;
soundness did not restrict $I$ to be even. These two details are recorded
explicitly in the inherited independent arithmetic audit.
\end{remark}

We can now finish completeness for both new equations. Starting with a
finite ZFC refutation, use the inherited canonical data
\eqref{eq:canonical}. Lemma~\ref{lem:joint-positive} selects $t$
with $P>0$. If $P=1$, choose $s=1,g=0$. If $P\ge2$, then
$3P-4>0$ is $2\pmod3$ and hence nonsquare; ordinary Pell existence
gives positive $s,g$ with
\[
 g^2-(3P-4)s^2=1.
\]
Apply Lemma~\ref{lem:interpolation} with $N=s$, then put
$r=I/(sc^2)\in\Z$. At this stage the pair $(I,f)$ and the coordinate
$r$ are fixed.

By intersectivity, choose $a_0$ with
$\pi(a_0+\gamma)\equiv0\pmod I$. The interpolation lemma supplies
$T\equiv a_0\pmod I$ and its integer auxiliary root $y$, without
changing $I,f,r$. The main canonical root $d$ satisfies
\[
 A+Bd=\Phi_H(HT+d-b_0)=H^5\pi(T+\gamma),
\]
so $(A+Bd)/S$ is an integer. Choose
\begin{equation}\label{eq:complete-final-coordinate}
 n=\frac{A+Bd}{S}+y+3f+9\tau+27g.
\end{equation}
One conjugate factor in \eqref{eq:primary} then vanishes. For
\eqref{eq:companion}, replace $y$ in
\eqref{eq:complete-final-coordinate} by $\delta y$. Every choice is
made from already fixed data: native coordinates, $t$, $s$, the pair
$(I,f)$ and $r$, the intersective residue, $(T,y)$, and finally $n$.
There is no circular dependence of the modulus on the requested residue.

Together with Section~\ref{sec:soundness}, this proves that each new
polynomial has an integer zero if and only if there is a finite
syntactic ZFC refutation, with the inherited ordinary calibration
identified above. The elementary arithmetic arguments and those
calibration proofs are formalizable in ZFC, yielding
\eqref{eq:target}. This does not supply a new machine-checked internal
ZFC derivation of the complete biconditional.

\section{The literal denominator-free norm}\label{sec:clearing}

We give a finite ring recipe for the denominator clearing in
\eqref{eq:primary}. This also specifies exactly what the executable
\texttt{norm5.py} computes. Use generic arguments
$x,a_1,a_2,d_0,e_0,N_0,\delta_0$, and work temporarily in the
free rank-four algebra with basis $1,u_0,v_0,u_0v_0$ and relations
$u_0^2=d_0$, $v_0^2=e_0$. For a four-tuple $z=(z_0,z_1,z_2,z_3)$,
its square has coordinates
\begin{align*}
 [z^2]_0&=z_0^2+d_0z_1^2+e_0z_2^2+d_0e_0z_3^2,\\
 [z^2]_1&=2(z_0z_1+e_0z_2z_3),\\
 [z^2]_2&=2(z_0z_2+d_0z_1z_3),\\
 [z^2]_3&=2(z_0z_3+z_1z_2).
\end{align*}
Let $\mathbf e=(1,0,0,0)$ and define
\begin{align*}
 \mathbf t&=(x^2+d_0+e_0,-2x,-2x,2),\\
 \mathbf q&=\delta_0\bigl(\mathbf t-(a_1+a_2)\mathbf e\bigr)
                      -N_0\mathbf e,\\
 r_0&=\delta_0^2a_1a_2,\qquad s_0=\delta_0N_0,\\
 \mathbf h&=\mathbf q^2-4\bigl(r_0+s_0(a_1+a_2)\bigr)\mathbf e,\\
 \mathbf f&=\mathbf h^2-64r_0s_0\mathbf t.
\end{align*}
If $\mathbf f=(f_0,f_1,f_2,f_3)$, put
\begin{align*}
 g_0&=f_0^2+e_0f_2^2-d_0(f_1^2+e_0f_3^2),\\
 g_1&=2(f_0f_2-d_0f_1f_3),\\
 \mathcal C(x;a_1,a_2,d_0,e_0,N_0,\delta_0)&=g_0^2-e_0g_1^2.
\end{align*}
Every operation in this recipe is addition, subtraction, or
multiplication in the coefficient ring. No root, division, or
data-dependent loop occurs in the exported arithmetic circuit.

The three-root identity
\[
 F_8(x;a,b,c)
 =\bigl((x^2-a-b-c)^2-4(ab+ac+bc)\bigr)^2-64abcx^2
\]
shows that $\mathbf f$ represents
$\delta_0^4F_8(x-u_0-v_0;a_1,a_2,N_0/\delta_0)$.
The final two quadratic norms therefore give the universal polynomial
identity
\begin{equation}\label{eq:cleared-norm-identity}
 \mathcal C
 =\delta_0^{16}\Fn(x;a_1,a_2,d_0,e_0,N_0/\delta_0).
\end{equation}
The right side is interpreted after clearing denominators, so the identity
is a polynomial identity even at $\delta_0=0$.
For the principal construction substitute
\[
 (x,a_1,a_2,d_0,e_0,N_0,\delta_0)
 =(Z,B^2\Dm,9S^2\Df,81S^2\Dt,729S^2\Dg,S^2\Cc,\delta).
\]
For the companion, replace the last two arguments by
$(S^2\delta\Cc,1)$.

The independent exact checker constructs $F_{32}$ by successive binomial
translations $F(x+\sqrt a)F(x-\sqrt a)$ and compares every integer
coefficient after denominator clearing. All $20349$ terms match
\eqref{eq:cleared-norm-identity}. Its intermediate term counts are
$2,6,35,495,20349$. This is an exact universal polynomial check;
the additional modular evaluations in the release serve as transcription
controls for the specialized front end.

\section{Attained degree and arithmetic cost}\label{sec:degree}

The degrees of the relevant registers are
\begin{center}\small
\begin{tabular}{@{}lrrrrrrrrrr@{}}\toprule
Register &$S$&$A$&$B$&$\Dm$&$\Df$&$\Dt$&$\Dg$&$\delta$&$V$&$\Cc$\\
\midrule
Degree &34&35&28&14&20&34&46&32&25&82\\\bottomrule
\end{tabular}
\end{center}
The norm $\Fn$ is even in $Z$ and homogeneous of degree $16$ in
$Z^2$ and the five radicands. Consequently its cleared version is a
homogeneous polynomial of degree $16$ in
\[
 \delta Z^2,\quad \delta B^2\Dm,\quad9\delta S^2\Df,
 \quad81\delta S^2\Dt,\quad729\delta S^2\Dg,\quad S^2\Cc.
\]
These six input degrees are, respectively,
\[
 102,\quad102,\quad120,\quad134,\quad146,\quad150.
\]
The coefficient of the sixteenth power of the last input is one.
Therefore $(S^2\Cc)^{16}$ is the entire degree-$2400$ contribution;
every other contribution has smaller degree. This yields the degree
upper bound and reduces attainment to a single leading-form check.

Write
\begin{align*}
 q_0&=6\beta j,& W_{00}&=12w,\\
 a_0&=2q_0^4W_{00}v
     =2^7 3^5\beta^4j^4wv,& I_0&=src^2,\\
 R_0&=-q_0^3(W_{00}C-q_0z),& b&=4I_0T-R_0.
\end{align*}
The leading forms of $\Delta,H,Q,\Df,V$ are respectively
$a_0^2,4a_0,a_0^4I_0^2,a_0^2I_0^2,a_0^2I_0^2b$.
It follows that the leading form of $\primary$ is exactly
\begin{equation}\label{eq:primary-leading}
 2^{320}a_0^{288}I_0^{128}b^{32}.
\end{equation}
Selecting the maximal $T$ power in $b^{32}$ gives the explicit coefficient
\begin{equation}\label{eq:primary-coefficient}
 \bigl[j^{1152}w^{288}v^{288}s^{160}r^{160}c^{320}T^{32}\bigr]
 \primary
 =2^{2400}3^{1440}\beta^{1152}>0.
\end{equation}
The exponents sum to $2400$. Since no other fixed port occurs in this
coefficient, specializing to the actual six-tuple cannot make it vanish.

For the companion, the radicand $S^2\delta\Cc$ has degree $182$,
uniquely above the other radicands and $Z^2$. Its sixteenth power is
the degree-$2912$ part, with leading form
\[
 2^{320}a_0^{352}I_0^{160}b^{32}.
\]
In particular,
\begin{equation}\label{eq:companion-coefficient}
 \bigl[j^{1408}w^{352}v^{352}s^{192}r^{192}c^{384}T^{32}\bigr]
 \companion
 =2^{2848}3^{1760}\beta^{1408}>0.
\end{equation}
The exponents sum to $2912$. These are attained degrees of the
specialized polynomials, not only syntactic degree estimates for circuits
with generic coefficient ports.

The final live-node census is
\begin{center}
\begin{tabular}{@{}lrrrr@{}}\toprule
Polynomial & Additions & Multiplications & Subtractions & Total\\
\midrule
$\primary$ &44&139&32&215\\
$\companion$ &44&133&32&209\\\bottomrule
\end{tabular}
\end{center}
All declared variables and all six coefficient leaves are live. The
counts include the complete denominator clearing and the complete norm,
not just the arithmetic front end. A multiplication by a fixed scalar
is counted, as is each square; already computed registers may be reused.
The exporter removes unused registers and coalesces identical operations.
No assertion of optimal circuit complexity is made.
